% character_breakdowns.tex — Character breakdowns for casting

\section{Character Breakdowns}

\subsection{Franz Kafka}
\textbf{Lead. Male. 29--40 across the film. Jewish. Czech-German.}

The gravitational center of the film. Kafka is not a conventional ``genius'' character. He is brilliant, dryly funny, profoundly ashamed, emotionally avoidant, needy in ways he cannot admit, and afraid of intimacy in direct proportion to how desperately he craves it. He is vain in small, embarrassing ways --- pleased by praise, mortified by his own pleasure. He is self-lacerating without self-pity. He is a man who understands systems better than anyone and cannot escape any of them.

Physically: thin, angular, enormous dark eyes. The kind of face that cameras love. The body is not decorative; it is the site of the film's central conflict. He deteriorates visibly across the film --- from lean and alert to gaunt and fragile to, in the final act, skeletal.

The performance must carry six major awards-clip scenes and numerous passages of minimal dialogue. This is a role built on micro-expression, silence, and the visible architecture of thought. The actor must make shame interesting to watch.

\textbf{Designed for: Adrien Brody.}

\subsection{Hermann Kafka}
\textbf{Major supporting. Male. 50s--60s. Jewish. Czech-German.}

Franz's father. A self-made businessman who runs a fancy goods shop in Prague. Large, loud, practical, energetically alive. He is not a villain. He is, in many ways, admirable: hardworking, charming with customers, generous in his fashion. His tragedy --- and his son's --- is that his virtues are themselves crushing. His practicality, his appetite for life, his confidence, his expectations --- all of these land on Franz not as abuse but as weather. Unavoidable, impersonal, and devastating.

Hermann must never be played as a monster. He must be played as a father who loves his son in the only language he has, and whose love is a language his son cannot speak.

\subsection{Julie Kafka}
\textbf{Supporting. Female. 50s. Jewish. Czech-German.}

Franz's mother. A diplomat without a treaty. She has spent her life managing the gap between Hermann and the world, and between Hermann and Franz. She is warm, tired, practical, and more clear-eyed than anyone in the family gives her credit for.

Her pivotal actions define the title engine of the film. First: refusing to deliver the Letter to His Father. This is not cowardice. It is the painful clarity of a woman who has spent thirty years absorbing the friction between two people and who understands that the truth, administered all at once in writing, will destroy rather than heal. ``If I carry this into that room, he will hear accusation before he hears pain,'' she tells Franz. She puts the letter in a drawer. Years later, unable to keep serving as the family's courier of pain, she mails the undelivered packet back to Berlin. Both moments must be played as choices, not inevitabilities.

\subsection{Ottla Kafka}
\textbf{Supporting. Female. Early 20s. Jewish. Czech-German.}

Franz's youngest sister and closest ally. Warm, direct, slightly rebellious. She wants to work on a farm, which horrifies Hermann. She understands Franz better than anyone in the family and loves him without trying to fix him --- but she is also the only person who will tell him the truth to his face. Her confrontation after the Letter (``You wrote a hundred pages because paper doesn't interrupt. Then you gave them to Mother because she does.'') is one of the film's most devastating moments, precisely because it comes from the person who loves him most. She is the family member who actually defied Hermann in person, which gives her moral authority her brother lacks.

\subsection{Max Brod}
\textbf{Major supporting. Male. Late 20s--30s. Jewish. Czech-German.}

Kafka's closest friend and literary executor. A writer himself --- prolific, sociable, optimistic. Everything Kafka is not on the surface, and devoted to him in a way that transcends professional admiration. Brod recognized Kafka's genius before anyone, including Kafka.

His central dramatic function is the manuscripts dilemma, but the film earns that dilemma through a genuine philosophical disagreement that develops across three movements. Brod believes literature belongs to the world; Kafka believes it belongs to the writer. Their final confrontation is not a simple request and capitulation --- it is a clash between two defensible positions about art, privacy, and what we owe the dead. Brod argues that destruction is its own vanity. Kafka argues that prayers are private. When Brod agrees to burn the manuscripts, both men may know he won't. The actor must make both sides of every decision equally visible.

\subsection{Felice Bauer}
\textbf{Supporting. Female. Mid 20s. Jewish. German.}

Kafka's first serious romantic partner. Practical, organized, intelligent. She works for a dictation machine company and runs the Berlin office. We must see her in command of that world --- machines, timetables, correspondence, adult competence --- because she represents social adulthood, the ordinary life Kafka most desires and most fears. Critically, she is not a passive recipient of Kafka's withdrawal. She reads his letters to a mutual friend, in which he describes the impossibility of the life she is offering, and she ends the engagement herself. Her exit line (``I hope you find what you're looking for. I think you hope you don't.'') is one of the most perceptive judgments in the film. She must be played as a woman of clarity and dignity, not a victim.

\subsection{Milena Jesenská}
\textbf{Supporting. Female. Mid 20s. Czech. Non-Jewish.}

A journalist, translator, and intellectual. She translates Kafka's work into Czech and, through the translations, sees into him with a clarity that is both thrilling and terrifying. Where Felice offers a concrete adult life, Milena offers mutual reading --- the terrifying ecstasy of being understood at the level of language, desire, and thought. She is passionate, direct, perceptive, unhappily married. Her scenes with Kafka have an electric, dangerous intimacy --- she is the person who comes closest to knowing him, and that proximity is precisely what drives him away.

\subsection{Dora Diamant}
\textbf{Supporting. Female. Mid 20s. Jewish. Polish.}

The last love. Warm, practical, young but not naive. From a Hasidic family, which she has partially left behind. She meets Kafka at a summer camp and moves with him to Berlin, where they experience the closest thing to domestic happiness Kafka has ever known. Where Felice is order and Milena is dangerous recognition, Dora is bodily care and present-tense life: food, warmth, Hebrew lessons, shared space, someone still there in the morning. She loves him without the tortured dynamics of his earlier relationships. She sees him clearly and stays anyway. The actor must convey both tenderness and the quiet devastation of watching someone you love die.
